\documentclass[11pt]{article}
\usepackage[margin=2.1cm]{geometry}
\usepackage{booktabs}
\usepackage{amsmath}
\usepackage{enumitem}
\usepackage[T1]{fontenc}
\usepackage{lmodern}

\setlength{\parindent}{0pt}
\setlength{\parskip}{5pt}
\newcommand{\code}[1]{\texttt{#1}}
\newcommand{\answerbox}[1]{\vspace{2pt}\fbox{\parbox[c][#1][c]{\dimexpr\linewidth-2\fboxsep-2\fboxrule}{\ }}\vspace{4pt}}

\begin{document}

\begin{center}
{\Large\bfseries EE\,362 --- Pre-Lab 10}\\[2pt]
{\large Instruction Decoder I}\\[6pt]
\begin{tabular}{@{}p{6.5cm}p{4cm}p{4cm}@{}}
Name: \hrulefill & Section: Mon / Tue & Date: \hrulefill
\end{tabular}
\end{center}

\vspace{2pt}\hrule\vspace{4pt}

{\small\textbf{Due at the start of the lab session.} Five points, about twenty
minutes, on paper. Bring the instruction-format sheet.}

\vspace{4pt}\hrule

\section*{Question 1 --- Pull the instruction apart}

\textbf{(a)} Split each instruction into its fields.

\begin{center}
\begin{tabular}{lllllll}
\toprule
word & format & opcode & rs & rt & rd / immediate & func / sa \\
\midrule
\code{0x00221820} & \rule{1.4cm}{0.4pt} & \rule{1.6cm}{0.4pt} & \rule{1cm}{0.4pt} & \rule{1cm}{0.4pt} & \rule{2cm}{0.4pt} & \rule{1.6cm}{0.4pt} \\[6pt]
\code{0x2021FFFF} & \rule{1.4cm}{0.4pt} & \rule{1.6cm}{0.4pt} & \rule{1cm}{0.4pt} & \rule{1cm}{0.4pt} & \rule{2cm}{0.4pt} & \rule{1.6cm}{0.4pt} \\[6pt]
\code{0x00021040} & \rule{1.4cm}{0.4pt} & \rule{1.6cm}{0.4pt} & \rule{1cm}{0.4pt} & \rule{1cm}{0.4pt} & \rule{2cm}{0.4pt} & \rule{1.6cm}{0.4pt} \\
\bottomrule
\end{tabular}
\end{center}

\textbf{(b)} Two of those write a register. \emph{Which field holds the
destination in each case}, and why is it not the same field for both?

\answerbox{2.0cm}

\textbf{(c)} The opcode always sits at bits 31--26, in every format. Give the
reason --- it is a hardware reason, not a convention.

\answerbox{1.6cm}

\section*{Question 2 --- The control table}

Fill in only the signals that DIFFER from the default. Write ``--'' for the
rest.

\begin{center}
\small
\begin{tabular}{lllllll}
\toprule
instr & ALU\_func & ALU\_Asel & IR\_Xsel & RF\_Dsel & RF\_strb & RF\_Wsel \\
\midrule
\code{ADD}  & \rule{1.5cm}{0.4pt} & \rule{1.5cm}{0.4pt} & \rule{1.3cm}{0.4pt} & \rule{1.3cm}{0.4pt} & \rule{1.5cm}{0.4pt} & \rule{1.3cm}{0.4pt} \\[6pt]
\code{ADDI} & \rule{1.5cm}{0.4pt} & \rule{1.5cm}{0.4pt} & \rule{1.3cm}{0.4pt} & \rule{1.3cm}{0.4pt} & \rule{1.5cm}{0.4pt} & \rule{1.3cm}{0.4pt} \\[6pt]
\code{LUI}  & \rule{1.5cm}{0.4pt} & \rule{1.5cm}{0.4pt} & \rule{1.3cm}{0.4pt} & \rule{1.3cm}{0.4pt} & \rule{1.5cm}{0.4pt} & \rule{1.3cm}{0.4pt} \\
\bottomrule
\end{tabular}
\end{center}

\textbf{(b)} \code{ADDI} and \code{LUI} differ on exactly \emph{one} of those
signals. Circle it, and say what would go wrong in the Lab 12 program if
\code{ADDI} were given \code{LUI}'s value.

\answerbox{2.2cm}

\section*{Question 3 --- The latch, and how to catch it}

A classmate writes their decoder without assigning the default first. They test
it by applying one instruction, checking the signals, applying the next, and so
on. Every test passes.

\textbf{(a)} Explain, in two or three sentences, what their hardware actually
does and why their testing cannot detect it.

\answerbox{2.4cm}

\textbf{(b)} Design the vector that \emph{would} catch it. Give the exact
sequence of instructions and the signal you would check.

\answerbox{2.2cm}

\textbf{(c)} You have met this bug twice before in this course. Name both labs
and say what was different about them.

\answerbox{2.0cm}

\vfill
{\small\emph{Question 3 is the one that matters.} It is the same question as
Section 5 of your lab report.}

\end{document}
